\documentclass[11pt,a4paper]{article}

\usepackage[margin=1in]{geometry}
\usepackage{enumitem}
\usepackage{hyperref}
\usepackage{titlesec}
\usepackage{xcolor}
\usepackage{tabularx}

% ----- NO INDENT -----
\setlength{\parindent}{0pt}
\setlength{\parskip}{4pt}

% ----- COLORS -----
\definecolor{accent}{RGB}{0, 102, 153}

\hypersetup{
    colorlinks=true,
    urlcolor=accent,
    linkcolor=accent
}

\setlist[itemize]{leftmargin=*, noitemsep, topsep=3pt}

% ----- SECTION STYLE -----
\titleformat{\section}
  {\large\bfseries\color{accent}}
  {}{0em}{}[\titlerule]

\titlespacing{\section}{0pt}{10pt}{6pt}

\begin{document}

% ===== HEADER =====
\begin{center}
    {\LARGE \textbf{Dominik Maszczyk}} \\[4pt]
    \textcolor{accent}{Senior Software Engineer} \\[6pt]
    Copenhagen, Denmark \\
    \href{https://www.linkedin.com/in/dominik-maszczyk/}{linkedin.com/in/dominik-maszczyk} \quad | \quad
    \href{https://github.com/Skitionek}{github.com/Skitionek}
\end{center}

\vspace{6pt}

% ===== PROFILE =====
\section*{Profile}
Senior engineer with 9+ years of experience building scientific data platforms at the intersection of IT and life science research.
I combine hands-on Python and full-stack development with technical leadership — contributing to roadmap, prioritization, and long-term sustainability in close collaboration with researchers.
I have worked closely with bio-chemical research groups, and I take pride in translating ambiguous scientific needs into maintainable solutions.

% ===== CORE COMPETENCIES =====
\section*{Core Competencies}
\begin{itemize}
    \item Technical leadership: contributing to roadmap and prioritization
    \item Scientific data platforms for genomics, proteomics, and bioinformatics research groups
    \item Multi-database architectures (PostgreSQL, Neo4j, ArangoDB, Redis, Elasticsearch)
    \item Distributed systems and message queues experience
    \item DevOps, CI/CD, and open-source tool integration
    \item Full-stack development (TypeScript, Angular, React, Python/Flask)
    \item Agile ways of working and cross-functional team collaboration
\end{itemize}

% ===== EXPERIENCE =====
\section*{Professional Experience}

\textbf{\href{https://www.biosustain.dtu.dk/}{Novo Nordisk Foundation Center for Biosustainability}} \hfill \textit{Oct 2016 -- Present}

\vspace{4pt}

\textit{Senior Software Engineer} \hfill \textit{Mar 2025 -- Present, Copenhagen}
\begin{itemize}
    \item Contributed to platform roadmap and prioritization in collaboration with researchers, bioinformaticians, and management.
    \item Designed systems for unified data storage, access control, and provenance tracking across scientific workflows.
    \item Improved reproducibility and traceability of research workflows through data architecture design.
    \item Maintained and scaled cloud-based research applications on Azure; introduced CI/CD pipelines.
\end{itemize}

\vspace{4pt}

\textit{Software Engineer} \hfill \textit{Jun 2022 -- Dec 2024, Copenhagen}
\begin{itemize}
    \item Core contributor to \href{https://lifelike.bio}{Lifelike platform} — a scientific knowledge graph and data visualization system for bioinformaticians; fork at \href{https://github.com/skitionek/lifelike}{skitionek/lifelike}.
    \item Authored \textbf{\~75\% of the codebase} across frontend, backend, and services (Python/Flask, Angular/TypeScript).
    \item Built data processing pipelines, modeling and generating aggregated graph databases (Ecocyc, Pubmed, KEGG etc.).
    \item Collaborated directly with scientists and bioinformaticians to translate research needs into platform features.
\end{itemize}

\vspace{4pt}

\textit{Software Engineer (On-site)} \hfill \textit{Feb 2022 -- May 2022, \href{https://ucsd.edu}{UC San Diego}}
\begin{itemize}
    \item Collaborated with distributed teams on platform development and integration.
\end{itemize}

\vspace{4pt}

\textit{Research Assistant → Software Engineer} \hfill \textit{Oct 2016 -- May 2022, Copenhagen}
\begin{itemize}
    \item Developed full-stack features for synthetic biology platform.
    \item Led infrastructure evolution:
    \begin{itemize}
        \item Docker Compose → Kubernetes migration
        \item DigitalOcean → Google Cloud transition
    \end{itemize}
    \item Maintained production systems supporting research workflows.
\end{itemize}

% ===== OPEN SOURCE =====
\section*{Open Source}

\textbf{\href{https://github.com/Skitionek/nf-mapper}{nf-mapper}} — Nextflow tooling
\begin{itemize}
    \item Improves structure and maintainability of workflow pipelines.
    \item Focused on developer experience and workflow introspection.
\end{itemize}

\vspace{4pt}

\textbf{More projects:} \href{https://github.com/Skitionek}{github.com/Skitionek}

% ===== PUBLICATIONS =====
\section*{Publications}
\begin{itemize}
    \item \textbf{\href{https://www.researchgate.net/publication/400970354_PREFER_An_Ontology_for_the_PREcision_FERmentation_Community}{PREFER: An Ontology for the PREcision FERmentation Community}}
    \item \textbf{\href{https://zenodo.org/records/10103768}{Reactome modified for tracing (ArangoDB version)}}
    \item \textbf{\href{https://github.com/biosustain/lifelike}{Lifelike v0.98.1}}
\end{itemize}

% ===== EDUCATION =====
\section*{Education}

\textbf{M.Sc. Electrical and Electronics Engineering (Double Degree)} \\
\href{https://www.kaist.ac.kr/en/}{KAIST} \& \href{https://www.dtu.dk/english}{Technical University of Denmark (DTU)} \hfill 2016 -- 2018

\vspace{4pt}

\textbf{B.Sc. Electrical Engineering} \\
\href{https://pg.edu.pl/en}{Gdańsk University of Technology} \& \href{https://www.kit.edu/english/}{Karlsruhe Institute of Technology} \hfill 2012 -- 2016

% ===== LANGUAGES =====
\section*{Languages}
English (Professional), Polish (Native), Danish (Beginner)

\end{document}